\subsection{模的投射维数}

定义 1. --- 设 M 是一个 A-模。称 M 的投射维数，并记作 dpA(M)，为 M 的投射分解长度在 Z 中的下确界（X，第 48 页）。

我们因此有 \(\mathrm{dp}_{\mathbf{A}}(0) = -\infty\) , \(\mathrm{dp}_{\mathbf{A}}(\mathbf{M}) \geqslant 0\) （若 \(\mathbf{M} \neq 0\) ）。为使 \(\mathbf{M}\) 是投射的，必要且充分条件是 \(\mathrm{dp}_{\mathbf{A}}(\mathbf{M}) \leqslant 0\) .

引理 1. --- 若 dp \(_{A}\) (M) \textless{} n \textless{} +∞，则对每个 A-模 N 有 Ext \(_{A}^{n}\) (M, N) = 0，且对每个右 A-模 P 有 Tor \(_{n}^{A}\) (P, M) = 0。

这由 M 具有长度为 \textless{} n 的投射分解以及 X，第 100 页，定理 1 直接得出。

命题 1. --- 设 M 是一个 A-模，n 为整数 ≥ 0。以下条件等价：

\begin{enumerate}
\def\labelenumi{(\roman{enumi})}
\item
  \(\mathrm{dp}_{\mathbf{A}}(\mathbf{M}) \leqslant n\) （即（定义 1），M 具有长度为 \(\leqslant n\) 的投射分解）
\item
  \(\operatorname{Ext}_{\mathbf{A}}^{r}(\mathbf{M},\mathbf{N}) = 0\) （对每个 A-模 N 及每个整数 \(r > n\) ）;
\item
  \(\operatorname{Ext}_{\mathbf{A}}^{n+1}(\mathbf{M}, \mathbf{N}) = 0\) （对每个 A-模 \(\mathbf{N}\) ）;
\item
  对每个正合序列
\end{enumerate}

\[
0 \rightarrow \mathrm{K} \rightarrow \mathrm{P} _ {n - 1} \rightarrow \dots \rightarrow \mathrm{P} _ {0} \rightarrow \mathrm{M} \rightarrow 0
\]

其中 \(P_{i}\) 是投射的，K 是投射的。

\begin{enumerate}
\def\labelenumi{(\roman{enumi})}
\item
  \(\Rightarrow\) (ii)：这由引理 1 得出。
\item
  \(\Rightarrow\) （iii）：这是平凡的。
\item
  \(\Rightarrow\) (iv)：在(iv)的情形下，对每个 A-模 N，有一个从 \(\operatorname{Ext}_{\mathbf{A}}^{1}(\mathbf{K},\mathbf{N})\) 到 \(\operatorname{Ext}_{\mathbf{A}}^{n + 1}(\mathbf{M},\mathbf{N})(\mathbf{X},\mathbf{p}.128,\operatorname{cor.}4)\) 上的同构；若满足(iii)，则有
\end{enumerate}

\[
\operatorname{Ext} _ {\mathbf {A}} ^ {1} (\mathbf {K}, \mathbf {N}) = 0
\]

对所有 N 成立，从而 K 是投射的（X，第 93 页，命题 10）。

\begin{enumerate}
\def\labelenumi{(\roman{enumi})}
\setcounter{enumi}{3}
\tightlist
\item
  \(\Rightarrow\) (i)：考虑正合序列（X，第 50 页）
\end{enumerate}

\[
0 \to Z _ {n - 1} (\mathrm{M}) \to \mathrm{L} _ {n - 1} (\mathrm{M}) \to \dots \to \mathrm{L} _ {0} (\mathrm{M}) \to \mathrm{M} \to 0.
\]

如果(iv)满足， \(Z_{n-1}(M)\) 是投射的，且 M 具有长度为 \(\leqslant n\) 的投射分解 .

推论 1. --- 设 \((\mathbf{M}_i)_{i\in \mathbb{E}}\) 为一族 A-模。我们有

\[
\mathrm{dp} _ {\mathbf {A}} \left(\bigoplus_ {i \in \mathbf {E}} \mathbf {M} _ {i}\right) = \sup _ {i \in \mathbf {E}} \mathrm{dp} _ {\mathbf {A}} (\mathbf {M} _ {i}).
\]

这由命题 1 的条件 (i) 与 (iii) 的等价性，以及 X，第 89 页，命题 7 得出。

在以下陈述中，我们一致认为 \(\pm\infty + 1 = \pm\infty - 1 = \pm\infty\) .

推论 2. --- 设

\[
0 \rightarrow \mathrm{M} ^ {\prime} \rightarrow \mathrm{M} \rightarrow \mathrm{M} ^ {\prime \prime} \rightarrow 0
\]

是A-模的正合序列。

\begin{enumerate}
\def\labelenumi{\alph{enumi})}
\tightlist
\item
  我们有
\end{enumerate}

\[
\mathrm{dp} _ {\mathbf {A}} (\mathbf {M}) \leqslant \sup \left(\mathrm{dp} _ {\mathbf {A}} (\mathbf {M} ^ {\prime}), \mathrm{dp} _ {\mathbf {A}} (\mathbf {M} ^ {\prime \prime})\right).
\]

一旦 \(\mathrm{dp}_{\mathbf{A}}(\mathbf{M}^{\prime \prime})\neq \mathrm{dp}_{\mathbf{A}}(\mathbf{M}^{\prime}) + 1\) ，等号便成立

\begin{enumerate}
\def\labelenumi{\alph{enumi})}
\setcounter{enumi}{1}
\tightlist
\item
  我们有
\end{enumerate}

\[
\mathrm{dp} _ {\mathbf {A}} \left(\mathrm{M} ^ {\prime \prime}\right) \leqslant \sup \left(\mathrm{dp} _ {\mathbf {A}} (\mathrm{M}), \mathrm{dp} _ {\mathbf {A}} \left(\mathrm{M} ^ {\prime}\right) + 1\right).
\]

一旦 \(\mathrm{dp}_{\mathbf{A}}(\mathbf{M}) \neq \mathrm{dp}_{\mathbf{A}}(\mathbf{M}')\) ，等号便成立 .

\[
\mathrm{dp} _ {\mathbf {A}} \left(\mathbf {M} ^ {\prime}\right) \leqslant \sup \left(\mathrm{dp} _ {\mathbf {A}} (\mathbf {M}), \mathrm{dp} _ {\mathbf {A}} \left(\mathbf {M} ^ {\prime \prime}\right) - 1\right)
\]

一旦 \(\mathrm{dp}_{\mathbf{A}}(\mathbf{M}) \neq \mathrm{dp}_{\mathbf{A}}(\mathbf{M}^{\prime \prime})\) ，等号便成立 .

我们来证明，例如 a) 部分，b) 和 c) 部分的证明是类似的。

\[
\geqslant 0.
\]

\[
\begin{array}{r l}\operatorname{Ext} _ {\mathbf {A}} ^ {n + 1} (\mathrm{M} ^ {\prime \prime}, \mathrm{N})&\rightarrow \operatorname{Ext} _ {\mathbf {A}} ^ {n + 1} (\mathrm{M}, \mathrm{N}) \rightarrow \operatorname{Ext} _ {\mathbf {A}} ^ {n + 1} (\mathrm{M} ^ {\prime}, \mathrm{N}) \rightarrow \operatorname{Ext} _ {\mathbf {A}} ^ {n + 2} (\mathrm{M} ^ {\prime \prime}, \mathrm{N})\\&\rightarrow \operatorname{Ext} _ {\mathbf {A}} ^ {n + 2} (\mathrm{M}, \mathrm{N}).\end{array}
\]

若 \(\mathrm{dp}_{\mathbf{A}}(\mathbf{M}')\) , \(\mathrm{dp}_{\mathbf{A}}(\mathbf{M}'')\leqslant n\) ，则 \(\operatorname{Ext}_{\mathbf{A}}^{n+1}(\mathbf{M}',\mathbf{N})=0\) 和 \(\operatorname{Ext}_{\mathbf{A}}^{n+1}(\mathbf{M}''\) , \(\mathbf{N})=0\) （命题 1），因此 \(\operatorname{Ext}_{\mathbf{A}}^{n+1}(\mathbf{M},\mathbf{N})=0\) 且 \(\mathrm{dp}_{\mathbf{A}}(\mathbf{M})\leqslant n\) （命题 1），因此

\[
\mathrm{dp} _ {\mathbf {A}} (\mathbf {M}) \leqslant \sup \left(\mathrm{dp} _ {\mathbf {A}} (\mathbf {M} ^ {\prime}), \mathrm{dp} _ {\mathbf {A}} (\mathbf {M} ^ {\prime \prime})\right).
\]

若 \(\mathrm{dp}_{\mathbf{A}}(\mathbf{M}) < \sup (\mathrm{dp}_{\mathbf{A}}(\mathbf{M}^{\prime}),\mathrm{dp}_{\mathbf{A}}(\mathbf{M}^{\prime \prime}))\) 则必然 \(\mathrm{dp}_{\mathbf{A}}(\mathbf{M}) < + \infty\) 。对每个 \(n > \mathrm{dp}_{\mathbf{A}}(\mathbf{M})\) 及每个 A-模 N，我们有

\[
\operatorname{Ext} _ {\mathbf {A}} ^ {n + 1} \left(\mathrm{M} ^ {\prime}, \mathrm{N}\right) \neq 0 \Leftrightarrow \operatorname{Ext} _ {\mathbf {A}} ^ {n + 2} \left(\mathrm{M} ^ {\prime \prime}, \mathrm{N}\right) \neq 0,
\]

由前面的正合序列；根据命题 1，这立即蕴含 \(\mathrm{dp}_{\mathbf{A}}(\mathbf{M}^{\prime \prime}) = \mathrm{dp}_{\mathbf{A}}(\mathbf{M}^{\prime}) + 1\) ，因为量 \(\mathrm{dp}_{\mathbf{A}}(\mathbf{M}^{\prime})\) , \(\mathrm{dp}_{\mathbf{A}}(\mathbf{M}^{\prime \prime})\) 之一 \(> \mathrm{dp}_{\mathbf{A}}(\mathbf{M})\) .

例. --- 设 \(a\) 是 \(\mathbf{A}\) 中既不可逆也不是右零因子的一个元素。那么 \(\mathrm{dp}_{\mathbf{A}}(\mathbf{A}/\mathbf{A}a)=1\) .

事实上，由正合序列 \(0 \to \mathbf{A}_s \xrightarrow{\varphi} \mathbf{A}_s \to \mathbf{A} / \mathbf{A}a \to 0\) （其中 \(\varphi(x) = xa\) ），我们有 \(\mathrm{dp}_{\mathbf{A}}(\mathbf{A} / \mathbf{A}a) \leqslant 1\) 。如果 \(\mathrm{dp}_{\mathbf{A}}(\mathbf{A} / \mathbf{A}a) < 1\) ，则 \(\mathbf{A} / \mathbf{A}a\) 是投射的，并且存在一个 A-线性映射 \(\psi : \mathbf{A}_s \to \mathbf{A}_s\) 使得 \(\psi \circ \varphi = 1d\) ；这意味着

\[
1 = \psi (\varphi (1)) = \psi (a) = a. \psi (1)
\]

且 a 是可逆的。

命题 2. --- 设 A 是左诺特的。设 M 是有限生成的 A-模，n 是整数 ≥ 0。以下条件等价：

\begin{enumerate}
\def\labelenumi{(\roman{enumi})}
\tightlist
\item
  \(\mathrm{dp}_{\mathbf{A}}(\mathbf{M}) \leqslant n\) .
\end{enumerate}

(i bis) M 有一个长度 ≤ n 的投射分解 P，使得对每个 i， \(P_{i}\) 都是有限生成的 A-模。

\begin{enumerate}
\def\labelenumi{(\roman{enumi})}
\setcounter{enumi}{1}
\item
  \(\operatorname{Ext}_{\mathbf{A}}^{r}(\mathbf{M},\mathbf{N}) = 0\) （对每个有限生成的 A-模 N 和每个整数 \(r > n\) ）.
\item
  \(\operatorname{Ext}_{\mathbf{A}}^{n+1}(\mathbf{M}, \mathbf{N}) = 0\) （对每个有限生成的 A-模 \(\mathbf{N}\) ）.
\item
  \(\operatorname{Tor}_r^{\mathbf{A}}(\mathbf{P},\mathbf{M}) = 0\) （对每个右 A-模 P 及每个 \(r > n\) ）.
\item
  \(\operatorname{Tor}_{n+1}^{\mathbf{A}}(\mathbf{A}/\mathfrak{a},\mathbf{M}) = 0\) （对每个有限生成的右理想 \(\mathfrak{a}\) （\(\mathbf{A}\) 的））.
\end{enumerate}

(i bis) \(\Rightarrow\) (i)：这是平凡的。

\begin{enumerate}
\def\labelenumi{(\roman{enumi})}
\item
  \(\Rightarrow\) (ii)：这由引理 1 得出。
\item
  \(\Rightarrow\) （iii）：这是平凡的。
\item
  \(\Rightarrow\) (i)：由(iii)及 X，第 107 页，命题 5，我们有 \(\operatorname{Ext}_{\mathbf{A}}^{n+1}(\mathbf{M}, \mathbf{N}) = 0\) （对每个 A-模 N），从而由命题 1 得(i)。
\item
  \(\Rightarrow\) (iv)：这由引理 1 得出。
\item
  \(\Rightarrow\) (v)：这是平凡的。
\item
  \(\Rightarrow\) (i bis)：设 \((\mathbf{L}, d)\) 是 \(\mathbf{M}\) 的一个自由分解，使得 \(\mathbf{L}_r\) 对每个 \(r\) 都是有限生成的（X，第 53 页，命题 6）。设 \(\mathbf{K} = \mathbf{Z}_{n-1}(\mathbf{L})\) ；则 \(\mathbf{K}\) 作为 \(\mathbf{L}_{n-1}\) 的子模是有限生成的，并且我们有一个正合序列
\end{enumerate}

\[
0 \rightarrow \mathrm{K} \rightarrow \mathrm{L} _ {n - 1} \rightarrow \mathrm{L} _ {n - 2} \rightarrow \dots \rightarrow \mathrm{L} _ {1} \rightarrow \mathrm{L} _ {0} \rightarrow \mathrm{M} \rightarrow 0.\tag{1}
\]

由(v)及 X，第 131 页，推论 3，我们有 \(\operatorname{Tor}_{1}^{\mathbf{A}}(\mathbf{A}/\mathfrak{a},\mathbf{K}) = 0\) （对 A 的每个有限生成右理想 \(\mathfrak{a}\) ）。根据 X，第 74 页，定理 2，A-模 K 是平坦的；由于它是有限生成的，因此是有限表示的（X，第 10 页，命题 5），所以它是投射的（X，第 13 页，推论），因此（1）是 M 的一个投射分解。

推论. --- 设 A 是左诺特的，且令 \(\mathcal{C}_{0}\) (相应地， \(\mathcal{C}\) ) 为有限生成投射 A-模（相应地，具有有限投射维数的有限生成 A-模）的类构成的集合。则 Grothendieck 群的同态 \(\mathrm{K}(\mathcal{C}_{0}) \to \mathrm{K}(\mathcal{C})\) 是双射。

这由 X，第 58 页，定理 1 得出（注意 \(\mathcal{C}_0\) 与 \(\mathcal{C}\) 由推论 2 可知是左正合的）。

